% AUTO-GENERATED by analyze_study2_secondary.py. DO NOT EDIT.
\begin{table*}[t]
\centering
\caption{Illustrative source-level contrast for
\studyclaim{StudyTwoCaseDesign}. The entries are the dominant SFT candidate and
dominant candidate from RF seed one for that design. Selection occurred after
manuscript assessment; this example explains a possible CAD mechanism but is not a
population or causal estimate.}
\label{tab:study2-rtl-case}
\small
\begin{tabularx}{\textwidth}{lXrrrrr}
\toprule
Arm & Observed RTL organization & Sample mass [\%] & Fmax [MHz] & LUT & FF & DSP \\
\midrule
SFT & Delayed samples and product registers followed by a parallel final sum &
\studyclaim{StudyTwoCaseSftMass} & \studyclaim{StudyTwoCaseSftFmax} &
\studyclaim{StudyTwoCaseSftLut} & \studyclaim{StudyTwoCaseSftFf} &
\studyclaim{StudyTwoCaseSftDsp} \\
RF seed one & Registered transposed multiply--accumulate chain &
\studyclaim{StudyTwoCaseRfMass} & \studyclaim{StudyTwoCaseRfFmax} &
\studyclaim{StudyTwoCaseRfLut} & \studyclaim{StudyTwoCaseRfFf} &
\studyclaim{StudyTwoCaseRfDsp} \\
\bottomrule
\end{tabularx}
\end{table*}
